\section{Introduction}
\label{sec:intro}

\textcolor{red}{
The rapid proliferation of diffusion-based generative models \cite{ho2020denoising,rombach2022high} has fundamentally transformed the landscape of digital content creation, enabling the synthesis of photorealistic images from textual prompts with unprecedented fidelity and diversity. Models such as Stable Diffusion, DALL-E, and Midjourney have democratized high-quality image generation, making it accessible to millions of users worldwide. However, this accessibility introduces pressing concerns regarding content provenance, copyright protection, and the potential for misuse through deepfakes and misinformation \cite{liu2025review, gupta2024comprehensive,bird2023typology,Rijsbosch2025watermarking}. As AI-generated imagery becomes perceptually indistinguishable from authentic photographs, the ability to reliably trace and verify the origin of digital content has emerged as a critical societal and technical challenge. Indeed, the urgency of this problem is underscored by recent regulatory developments, including the European Union's Artificial Intelligence Act (2024), which mandates that providers of generative AI systems mark their outputs in a machine-readable and detectable format \cite{euaiact2024}, which directs agencies to establish watermarking guidelines for AI-generated content.}

\textcolor{red}{
Digital watermarking has long served as the principal mechanism for asserting ownership and authenticity in multimedia content ~\cite{barni2001improved,cox2008digital}. Classical approaches embed imperceptible signatures into images by manipulating frequency-domain coefficients, most commonly through the Discrete Wavelet Transform (DWT) ~\cite{raval2003discrete,tian2023multi} or the Discrete Cosine Transform (DCT)~\cite{shensa2002discrete}, and verify their presence through handcrafted correlation or distance metrics. More recently, deep learning has given rise to end-to-end watermarking frameworks, exemplified by HiDDeN~\cite{zhu2018hidden}, which jointly train encoder–decoder architectures with differentiable noise layers to achieve improved robustness against common image distortions. More fundamentally, these post-hoc watermarking methods are applied \emph{after} image generation, rendering them detachable from the generative process and vulnerable to adversarial removal or jailbreaking~\cite{zhao2024invisible}.}

\textcolor{red}{
 In the context of generative models, a particularly promising paradigm has emerged: in-generation watermarking, wherein the watermark is embedded directly into the generative process rather than applied as a post-hoc modification. This strategy preserves the natural output distribution of the model while ensuring that the watermark propagates through the entire generation pipeline. Recent advances have sought to integrate watermarking directly into the diffusion generation pipeline. Tree-Ring Watermarking~\cite{wen2023tree_ring} embeds concentric ring patterns into the Fourier space of the initial noise vector, producing watermarks that are invisible and structurally invariant to several geometric transformations. The Stable Signature~\cite{fernandez2023stable} fine-tunes the latent decoder of diffusion models to embed a recoverable binary signature into all generated outputs. More recently, methods such as ZoDiac~\cite{zhang2024attack}, Gaussian Shading~\cite{yang2024gaussian}, and ROBIN~\cite{huang2024robin} have introduced further innovations in latent-space watermark injection, adversarial concealment, and distribution-preserving embedding.
}

\textcolor{red}{
Despite these advances, existing methods exhibit important limitations that hinder their practical deployment. First, approaches that rely on handcrafted detection metrics, such as frequency-domain correlation or distance thresholds, suffer from threshold instability: the optimal decision threshold varies substantially across different image transformations (e.g., JPEG compression, geometric warping, cropping), rendering a single fixed threshold unreliable for real-world verification. Second, most in-generation watermarking methods employ isotropic frequency structures (i.e. concentric rings, uniform grids, or random masks) that distribute watermark energy uniformly by radius. Such patterns are inherently vulnerable to directional distortions, including motion blur, anisotropic resampling, and compression artefacts that preferentially attenuate specific orientations. Third, the verification strategies adopted by current methods typically operate in a binary, threshold-dependent manner without producing calibrated confidence scores, limiting their utility in deployment scenarios that require probabilistic decision-making under heterogeneous and unknown distortion conditions.}

\textcolor{red}{
In this paper, we propose SpiderMark, a novel watermarking framework for diffusion-based image generation that addresses these limitations through a hybrid approach coupling algorithmic watermark injection with learning-based verification. SpiderMark operates entirely at the noise initialisation stage of the diffusion process and requires no modification to pretrained model weights, ensuring full compatibility with existing diffusion architectures and sampling procedures. The framework comprises three tightly integrated components. First, a structured spider-web watermark pattern (consisting of radial spokes and concentric polygonal rings) is injected into the Fourier representation of the initial latent noise prior to the denoising process. Unlike conventional ring-based patterns, the spider-web design distributes watermark energy across both radial and angular orientations, yielding inherent robustness to directionally biased degradations. Second, a paired verification dataset is constructed by generating matched watermarked and clean images under identical prompts and noise seeds, from which frequency-domain features are extracted for training. Third, a lightweight learned verifier based on a fine-tuned ResNet-18 backbone operates in the frequency domain, augmented with auxiliary similarity cues (PSNR and L1 distance computed against the known watermark pattern), to produce calibrated probabilistic confidence scores for watermark detection.}

\textcolor{red}{
The main contributions of this work are as follows:}

\begin{itemize}
    \item \textcolor{red}{\textbf{A novel hybrid watermarking framework} that unifies algorithmic frequency-domain injection with learning-based probabilistic verification into a single, model-agnostic pipeline. SpiderMark requires no modifications to pretrained diffusion model weights and is compatible with standard sampling procedures, making it readily deployable for existing generative systems.}

    \item \textcolor{red}{\textbf{The spider-web watermark pattern}, a geometrically structured frequency-domain design that introduces both radial and angular components. This interconnected lattice of spokes and rings provides directional robustness, spatial redundancy, and balanced frequency coverage, outperforming conventional ring-based and grid-based patterns under anisotropic and localised distortions.}

    \item \textcolor{red}{\textbf{A learned frequency-domain verifier} based on a fine-tuned ResNet-18 backbone that fuses deep spectral features with explicit similarity cues (PSNR and L1 distance) computed against the known watermark pattern. The verifier operates on the Fourier spectrum of the latent noise recovered via inverse DDIM, outputting a calibrated confidence score that enables stable watermark detection under a fixed decision threshold, without requiring access to prompts, noise seeds, or the diffusion generator at inference time.}
\end{itemize}

\textcolor{red}{
These contributions establish SpiderMark as a practical and generalisable solution for content attribution in diffusion-based generative models, addressing key limitations of existing approaches with respect to robustness, verification reliability, and cross-dataset transferability. The rest of this paper is organized as follows. Section 2 reviews related work. Section 3 details the SpiderMark methodology, including watermark injection, spider-web pattern design, and the learned verification pipeline. Section 4 presents experimental results and ablation studies. Section 5 discusses implications for robustness evaluation and deployment. Section 6 concludes the paper and outlines future directions. }
